%% ma: L10-B14-0027
%% lop: 10
%% chuong: 5
%% bai: 14
%% dang: 10.14.5
%% loai: DS
%% muc: [TH, TH, NB, TH]
%% dap_an: "ĐSSĐ"
%% nguon: "chuyen-de-cac-so-dac-trung-cua-mau-so-lieu-khong-ghep-nhom-toan-10-chuong-trinh-moi.pdf (GV. Phan Nhật Linh) (P04), Bài 03 Dạng 1 Phần II Câu 4"
%% kien_thuc: "Lớp 10 Bài 13 (số trung bình, tứ phân vị), Bài 14 (khoảng biến thiên, khoảng tứ phân vị, giá trị bất thường)"
%% trang_thai: chinh-thuc
%% ghi_chu: "Tự duyệt theo yêu cầu giáo viên 22/10/2026."
\begin{ex}
Cuối học kì I, bạn An đạt kết quả sáu môn học như sau:
\begin{center}
\begin{tabular}{|c|c|c|c|c|c|c|}\hline
Môn & Toán & Văn & Anh & Lí & Hóa & Sinh\\\hline
Điểm TB & $7{,}2$ & $8{,}0$ & $5{,}8$ & $7{,}2$ & $9{,}0$ & $4{,}6$\\\hline
\end{tabular}
\end{center}
Xét tính đúng sai của các khẳng định sau:
\choiceTF{\True Điểm trung bình các môn học kì của bạn An, làm tròn đến hàng phần mười, là $7{,}0$}{Điểm trung bình các môn học kì của bạn An, làm tròn đến hàng phần mười, là $7{,}3$}{Khoảng biến thiên của bảng điểm của bạn An bằng $3{,}4$}{\True Khoảng tứ phân vị của bảng điểm của bạn An bằng $2{,}2$}
\loigiai{a) Đúng: $\overline x=\dfrac{41{,}8}6\approx6{,}97\approx7{,}0$.

b) Sai (theo ý a).

c) Sai: $R=9{,}0-4{,}6=4{,}4$.

d) Đúng: sắp xếp $4{,}6$; $5{,}8$; $7{,}2$; $7{,}2$; $8{,}0$; $9{,}0$; $Q_1=5{,}8$; $Q_3=8{,}0$; $\Delta_Q=2{,}2$.}
\end{ex}
